\section{Lecture Preprocessing and Hybrid Chunking}

\subsection{Speech Transcription as a Temporal Semantic Signal}

The first step is to convert lecture audio into a timestamped transcript. This transcript serves two purposes. At the semantic level, it provides the linguistic content needed for knowledge extraction and question answering. At the temporal level, it defines sentence boundaries and pauses that reveal how the speaker organizes the explanation. Unlike plain subtitle text, timestamped transcription preserves alignment between language and video time, which is essential for later evidence grounding.

For lecture understanding, the transcript is not used as an end product but as an intermediate representation. It helps determine where an explanation begins and ends, which terms are emphasized in a section, and how retrieved answers can be linked back to the lecture timeline.

In addition, sentence-level timestamps make it possible to aggregate transcript content into larger semantic units without losing temporal grounding. Once chunk boundaries are identified, the system can merge sentences into chunk-local contexts, which are later used for concept extraction and retrieval. This property is essential for lecture analysis because meaning is often distributed across several successive sentences rather than contained in a single utterance.

\subsection{Slide Transition Detection with Structural Similarity}

Visual structure is captured through slide transition detection. Let $I_t$ and $I_{t+1}$ denote two neighboring sampled frames. Their structural similarity is measured by the SSIM function:

\begin{equation}
\mathrm{SSIM}(I_t, I_{t+1}) =
\frac{(2\mu_t\mu_{t+1} + c_1)(2\sigma_{t,t+1} + c_2)}
{(\mu_t^2 + \mu_{t+1}^2 + c_1)(\sigma_t^2 + \sigma_{t+1}^2 + c_2)},
\end{equation}

where $\mu$ and $\sigma$ represent local means and variances, and $\sigma_{t,t+1}$ is the cross-covariance. A substantial drop in SSIM indicates that the frame structure has changed, which is often caused by a slide transition.

This criterion is well suited to lecture videos because it responds to structural changes rather than raw pixel differences. Minor pointer movement or speaker motion typically leaves the main slide layout unchanged, while a new slide produces a clear structural shift. As a result, SSIM provides a more robust transition signal than naive frame differencing.

The slide signal also carries an implicit discourse prior: when an instructor advances slides, the lecture usually enters a new subtopic, example, or proof step. Although this prior is imperfect, it provides a strong structural anchor that is especially useful when the transcript alone is ambiguous.

\subsection{Role of Slide Text and OCR}

Some lecture content appears visually but is only partially described in speech. Examples include equations, bullet-point lists, code-like notation, and slide titles that frame the current topic. For this reason, slide images can be processed into textual descriptions through OCR or multimodal interpretation, producing a complementary textual channel. Even when the report does not emphasize implementation details, this design choice has algorithmic importance: it reduces the gap between what is visible on the slide and what is spoken by the lecturer.

This additional signal becomes especially valuable in downstream retrieval. A student may ask about a term that appeared on a slide but was paraphrased verbally. By preserving slide text as an auxiliary evidence source, the system increases the chance of retrieving the relevant section.

\subsection{Hybrid Chunking}

Slide transitions alone do not define all meaningful boundaries, because a single slide may contain several conceptual stages, while silence-based segmentation alone may overreact to speaking rhythm and create fragments without clear semantic identity. LectureMind therefore combines complementary cues so that the resulting segments are both structurally plausible and semantically useful. We model chunking as boundary selection over candidates derived from visual and speech cues. Let $B_s$ denote slide-based candidates and $B_p$ denote pause-based candidates. A candidate boundary $b$ is scored by

\begin{equation}
S(b) = \alpha \, V(b) + \beta \, P(b) + \gamma \, C(b),
\end{equation}

where $V(b)$ measures visual evidence from slide change, $P(b)$ measures speech evidence such as pause duration or sentence completion, and $C(b)$ optionally reflects semantic coherence between neighboring regions. In practice, strong slide changes define high-confidence boundaries, large speech pauses add additional candidates within long slide intervals, and the final boundaries are filtered using coherence and duration constraints. This hybrid strategy preserves major lecture structure while recovering finer internal organization from speech rhythm, producing chunks that are concise enough for retrieval yet broad enough to capture a coherent teaching unit.

\begin{table}[H]
    \centering
    \caption{Comparison of boundary cues used in chunking}
    \begin{tabular}{@{}lll@{}}
        \toprule
        \textbf{Cue Type} & \textbf{Strength} & \textbf{Weakness} \\ \midrule
        Slide transition & Strong structural anchor & Misses within-slide topic shifts \\
        Speech pause & Captures local discourse rhythm & Sensitive to speaking style \\
        Semantic coherence & Encourages topic consistency & Depends on transcript quality \\ \bottomrule
    \end{tabular}
\end{table}